\section{Predicting Retrieval Failure}
\label{sec:prediction}

Section~\ref{sec:functional} measured retrieval under damage, and
Section~\ref{sec:theory} derived what retrieval should be if damage is blind.
This section asks three practical questions. How much retrieval survives, and
can it be predicted without simulation? Which individual queries fail? What
changes when the damage is aimed at hubs? A fourth question, raised at our
presentation, runs through all three: how certain are these numbers?

\subsection{How much retrieval survives}
\label{sec:prediction-aggregate}

Equation~\eqref{eq:r2} predicts the two-hop success of both query groups from
their redundancy distributions alone (Figure~\ref{fig:d-theory-sim}). Averaged
over 5,000 fresh random damage orders, the measured success rate differs from
the formula by at most 0.0009 for edge removal and 0.0010 for node removal, at
every tested damage level (up to $f=0.99$ for edges and $f=0.90$ for nodes).

\begin{figure}[!htbp]
  \centering
  \includegraphics[width=0.92\linewidth]{act2/primekg_5000_bfs_v1/fig_R_theory_vs_sim.pdf}
  \caption{Two-hop retrieval under random edge removal (solid; mean and one
  sample SD over the ten official seeds) and the zero-parameter formula
  $1-G_m(1-p^2)$ for each query group (dashed).}
  \label{fig:d-theory-sim}
\end{figure}

The same formula also predicts the restoration results of
Section~\ref{sec:restoration}, which were obtained independently. Random repair
of a share $b$ of the deleted edges simply raises the retention probability.
Relation-type repair spends the whole budget on deleted drug--protein and
disease--protein edges, the only edges a two-hop route uses. Substituting the
resulting retention probabilities into Eq.~\eqref{eq:r2} for the 60 test queries
predicts the gains in Table~\ref{tab:d-restoration}. All twelve
damage--budget--policy combinations fall within one seed SD of E's measurements.
The formula also explains why the simplest informed policy wins at $L=2$:
protein--protein repairs cannot change $R_2$, so any budget spent on them is
wasted for this task.

\begin{table}[!htbp]
\centering
\small
\caption{Two-hop restoration gains on held-out queries (percentage points):
zero-parameter prediction from Eq.~\eqref{eq:r2} against E's measurement
(mean $\pm$ SD over ten damage seeds).}
\label{tab:d-restoration}
\begin{tabular}{@{}ccrrrr@{}}
\toprule
& & \multicolumn{2}{c}{Random repair} & \multicolumn{2}{c}{Relation-type repair} \\
\cmidrule(lr){3-4}\cmidrule(lr){5-6}
Damage $f$ & Budget & Predicted & Measured & Predicted & Measured \\
\midrule
0.30 & 5\%  & 1.9  & $1.2 \pm 2.2$  & 7.9  & $7.0 \pm 3.9$ \\
0.30 & 10\% & 3.7  & $2.7 \pm 3.0$  & 15.6 & $13.0 \pm 6.2$ \\
0.30 & 20\% & 7.4  & $6.0 \pm 3.9$  & 30.2 & $27.5 \pm 7.1$ \\
0.50 & 5\%  & 3.0  & $3.0 \pm 2.6$  & 13.1 & $12.7 \pm 4.4$ \\
0.50 & 10\% & 6.1  & $6.0 \pm 4.0$  & 26.3 & $25.5 \pm 4.6$ \\
0.50 & 20\% & 12.3 & $13.2 \pm 2.8$ & 51.6 & $48.5 \pm 6.7$ \\
\bottomrule
\end{tabular}
\end{table}

\subsection{How certain are these numbers?}
\label{sec:prediction-uncertainty}

Our figures show wide shaded bands, yet we report agreement within 0.001. The
two statements answer different questions, and the report uses three different
measures of uncertainty. Table~\ref{tab:d-uncertainty} lists them for two-hop
retrieval under random edge removal.

\begin{table}[!htbp]
\centering
\small
\caption{Three measures of uncertainty for $R_2$ under random edge removal.
``Independent'' is the SD expected if the 200 queries failed independently with
their formula probabilities.}
\label{tab:d-uncertainty}
\begin{tabular}{@{}crrrrrr@{}}
\toprule
& & \multicolumn{2}{c}{One damage state} & & \multicolumn{2}{c}{Mean of estimate} \\
\cmidrule(lr){3-4}\cmidrule(lr){6-7}
$f$ & Formula & SD & Independent & Variance ratio & Ten seeds & 5,000 orders \\
\midrule
0.30 & 0.6631 & 0.046 & 0.031 & 2.2 & $0.680 \pm 0.012$ & $0.6628 \pm 0.0006$ \\
0.50 & 0.4134 & 0.047 & 0.032 & 2.2 & $0.433 \pm 0.013$ & $0.4140 \pm 0.0007$ \\
0.70 & 0.1796 & 0.035 & 0.025 & 1.9 & $0.189 \pm 0.010$ & $0.1804 \pm 0.0005$ \\
\bottomrule
\end{tabular}
\end{table}

\begin{enumerate}[leftmargin=*,itemsep=3pt]
\item \textbf{The spread of one damage state} is what the shaded bands show. It
answers the question: if this graph is damaged once, how far can the result
fall from the average? At $f=0.5$ the SD is 0.047, which is large because the
benchmark is small. With 200 independent queries the binomial SD would already
be 0.032. The variance is 2.2 times larger still, because queries share edges.
The 200 queries use only 39 drugs, an average of 5.1 queries per drug, so losing
one drug--protein edge can break several queries at once. The benchmark
therefore behaves like about 90 independent queries. This spread is a property
of the task, and it does not measure model error.
\item \textbf{The standard error of the ten-seed mean} is about 0.013. The
official ten-seed curves lie 1.4 to 2.5 percentage points above the formula for
$0.1\le f\le0.6$, which is 1.3 to 2.3 standard errors. The offset has the same
sign at every level because the damage states are nested: each seed removes a
prefix of one fixed order, so a seed that is lucky at $f=0.3$ tends to stay
lucky at $f=0.5$. The seed deviations at these two levels have correlation
0.58. Ten nested seeds are effectively fewer than ten independent samples.
\item \textbf{The expectation} is estimated from 5,000 fresh, independent
orders. Its standard error is below 0.0007, and it matches the formula within
0.001.
\end{enumerate}

Why should the agreement be this tight? Equation~\eqref{eq:r2} is not a fitted
model. Given blind damage and routes that share no edge, it is an exact
statement about the expectation. The 5,000-order test therefore checks two
things: that our code implements the retrieval rule correctly, and that its
assumptions hold on this graph. Both are verifiable ($\texttt{pairs\_with\_multiple\_relations}=0$,
and the deletion orders do not depend on the graph). Agreement at the level of
the Monte Carlo error is expected. A visible disagreement would have meant a
bug. The informative tests are the cases in which an assumption is violated
(Section~\ref{sec:prediction-attack}).

The same reasoning explains the numbers that differed between our slides. The
hop scan in Section~\ref{sec:functional} used three seeds. At $f=0.5$, $R_3$ is
84.5\% for those three seeds, 81.9\% for all ten, and 79.2\% over 600 fresh
orders. Wherever a long-run value exists, we treat it as the estimate, and we
treat small-seed values as single realisations.

\subsection{Which query fails?}
\label{sec:prediction-perquery}

The formula gives each query a survival probability $1-(1-p^2)^m$. We tested
it as a classifier on the official seeds 7--9 at 30\% and 50\% edge removal.
For comparison we trained a logistic regression on seeds 0--6, using the
logarithms of $m$, the mean and minimum degree of the shared proteins, and the
drug and disease degrees (Table~\ref{tab:d-perquery}).

\begin{table}[!htbp]
\centering
\small
\caption{Per-query prediction of two-hop failure under random edge removal
(test seeds 7--9, 600 query outcomes per level). The ceiling is the expected AUC
if outcomes were independent draws from the formula probabilities.}
\label{tab:d-perquery}
\begin{tabular}{@{}crrrrr@{}}
\toprule
Damage $f$ & Formula AUC & Logistic AUC & Ceiling & Formula Brier & Logistic Brier \\
\midrule
0.30 & 0.713 & 0.712 & 0.723 & 0.186 & 0.186 \\
0.50 & 0.720 & 0.712 & 0.712 & 0.204 & 0.204 \\
\bottomrule
\end{tabular}
\end{table}

An AUC near 0.72 may look modest, but it is the ceiling. Under blind damage, two
queries with the same $m$ have the same survival probability, and which of them
fails is decided by chance. No feature of the intact graph can predict that.
Both models reach the ceiling, and the five degree features add nothing to the
single number $m$. For random damage, the redundancy count is the complete
predictor.

\subsection{Targeted attacks break independence}
\label{sec:prediction-attack}

The node formula, Eq.~\eqref{eq:r2-node}, assumes blind damage. A degree
attack is not blind. After removing the top 5\% of nodes by degree, the formula
predicts that 87.9\% of queries survive, but 30.5\% do. We call the difference
the correlation penalty. It is 57.4 percentage points for the degree attack and
63.9 points for the betweenness attack.

Since $R_2^{\mathrm{node}}$ is the product of an endpoint factor $\phi^2$ and a
path factor, the penalty splits into two channels. A sequential split depends
on which factor is changed first, so we report the average of both orders
(Figure~\ref{fig:d-channels}). For the 5\% degree attack, lost endpoints account
for 17.4 points and lost paths for 40.0 points. Depending on the order, the
endpoint share ranges from 11 to 24 points and the path share from 34 to 46.
\begin{itemize}[leftmargin=*,itemsep=2pt]
\item \textbf{Endpoints.} Hub diseases are among the first nodes removed. The
disease endpoint of 27.5\% of queries is in the top 5\% by degree, compared with
7.5\% for drugs.
\item \textbf{Paths.} A shared protein must connect to both a drug and a
disease, which strongly favours hubs. The median degree of a query's shared
protein is 188 per path, against 25 for all proteins. A 5\% attack removes the
only shared protein of 84\% of single-path queries. It also fails 53\% of
multi-path queries: 36\% lose an endpoint, and 17\% keep both endpoints but
lose every route.
\end{itemize}
Path count is therefore not effective redundancy. Routes that pass through the
same class of hubs fail together, which is exactly the independence that
Eq.~\eqref{eq:r2} assumes. This is the retrieval counterpart of the structural
result in Section~\ref{sec:theory-null}: in both cases the graph hangs what
matters on its hubs.

\begin{figure}[!htbp]
  \centering
  \includegraphics[width=0.55\linewidth]{act2/primekg_5000_bfs_v1/fig_corr_channels.pdf}
  \caption{Correlation penalty, the blind-damage formula minus measured two-hop
  success, under targeted node removal, split into lost endpoints and lost
  paths. Each split averages the two orders of decomposition.}
  \label{fig:d-channels}
\end{figure}

Under attack, degree features become informative. In 20 repeats of 5-fold
cross-validation over queries, logistic regression reaches AUC 0.88 at the 5\%
degree attack, compared with 0.66 for the formula. It reaches 0.83 when folds
are split by drug, so that no drug appears in both training and test sets. We
treat this result as exploratory, because the features include the degree that
defines the attack.

\subsection{Structural metrics understate functional damage}
\label{sec:prediction-proxy}

Can the giant component stand in for retrieval when judging damage? At equal
damage it gets the direction right but the size wrong. At 5\% node removal, the
degree attack lowers $S$ by 6.6 percentage points more than random removal,
but it lowers $R_2$ by 56.5 points more. The functional effect is 8.5 times
larger. Both metrics agree that targeted attacks are worse than random
removal. Beyond 10\% damage the two attacks differ by at most two of the 200
queries, so neither metric can reliably rank them. To estimate retrieval
damage, the retrieval formula is the better tool. It needs only the redundancy
distribution of the benchmark, no simulation, and it is exact for blind damage.
Its failure under attack is itself a measurement: the size of the correlation
penalty shows how much the benchmark's redundancy depends on hubs.
